% Source-derived chapter 05: A lower bound on the number of degree-$n$ number fields that are monogenic / have a short generating vector
% Original source: squarefree-values-polynomial-discriminants
\chapter{A lower bound on the number of degree-$n$ number fields that are monogenic / have a short generating vector}
\label{squarefree-values-polynomial-discriminants-chapter-05}
\source{squarefree-values-polynomial-discriminants}

\begin{remark}[Source section]
\label{squarefree-values-polynomial-discriminants-chapter-05-source}
\source{squarefree-values-polynomial-discriminants}
\source{squarefree-values-polynomial-discriminants}
This chapter reproduces the corresponding section of the retrieved
LaTeX source, including its definitions, statements, examples, and
proofs. It has not yet been translated into Lean.
\notready
\end{remark}

% The source section title was: A lower bound on the number of degree-$n$ number fields that are monogenic / have a short generating vector
\label{latticearg}
\source{squarefree-values-polynomial-discriminants}

Let $g\in \Vmn(\R)$ be a monic real polynomial of degree $n$ and
nonzero discriminant with $r$ real roots and $2s$ complex roots.  Then
$\R[x]/(g(x))$ is naturally isomorphic to $\R^n\cong \R^r\times \C^s$
as $\R$-vector spaces via its real and complex embeddings (where we
view $\C$ as $\R+\R\sqrt{-1}$).  The $\R$-vector space $\R[x]/(g(x))$
also comes equipped with a natural basis, namely
$1,\theta,\theta^2,\ldots,\theta^{n-1}$, where $\theta$ denotes the
image of $x$ in $\R[x]/(g(x))$. Let $R_g$ denote the lattice spanned
by $1,\theta,\ldots,\theta^{n-1}.$ In the case that $g$ is an integral
polynomial in $\Vmn(\Z)$, the lattice $R_g$ may be identified with the
ring $\Z[x]/(g(x))\subset\R[x]/(g(x))\subset \R^n$.

Since $g(x)$ gives a lattice in $\R^n$ in this way, we may ask whether
this basis is reduced in the sense of Minkowski, with respect to the
usual inner product on $\R^n$.\footnote{Recall that a $\Z$-basis
  $\alpha_1,\ldots,\alpha_n$ of a lattice $L$ is called {\it
    Minkowski-reduced} if successively for $i=1,\ldots, n$ the vector
  $\alpha_i$ is the shortest vector in $L$ such that
  $\alpha_1,\ldots,\alpha_i$ can be extended to a $\Z$-basis of
  $L$. Most lattices have a unique Minkowski-reduced basis.}  More
generally, for any monic real polynomial $g(x)$ of degree $n$ and
nonzero discriminant, we may ask whether the basis
$1,\theta,\theta^2,\ldots,\theta^{n-1}$ is Minkowski-reduced for the
lattice $R_g$, up to a unipotent upper-triangular transformation
over~$\Z$ (i.e., when the basis $[1\;\;\theta\;\;\theta^2\;\cdots\;
  \theta^{n-1}]$ is replaced by $[1\;\;\theta\;\;\theta^2\;\cdots\;
  \theta^{n-1}]A$ for some upper triangular $n\times n$ integer matrix
$A$ with $1$'s on the diagonal).

More precisely, given $g\in \Vmn(\R)$ of nonzero discriminant, let us
say that the corresponding basis
$1,\theta,\theta^2,\ldots,\theta^{n-1}$ of $\R^n$ is {\it
  quasi-reduced} if there exist monic integer polynomials $h_i$ of
degree~$i$, for $i=1,\ldots,n-1$, such that the basis
$1,h_1(\theta),h_2(\theta),\ldots,h_{n-1}(\theta)$ of $R_g$ is
Minkowski-reduced (so that the basis
$1,\theta,\theta^2,\ldots,\theta^{n-1}$ is Minkowski-reduced up to a unipotent
upper-triangular transformation over~$\Z$). By abuse of language, we
then call the polynomial $g$ {\it quasi-reduced} as well. We say that
$g$ is {\it strongly quasi-reduced} if in addition $\Z[x]/(g(x))$ has
a unique Minkowski-reduced basis.

The relevance of being strongly quasi-reduced is contained in the
following lemma.

\begin{lemma}
\source{squarefree-values-polynomial-discriminants}
\notready\label{reducedg}
\source{squarefree-values-polynomial-discriminants}
Let $g(x)$ and $g^*(x)$ be distinct monic integer polynomials of
degree $n$ and nonzero discriminant that are strongly quasi-reduced
and whose $x^{n-1}$-coefficients vanish.  Then $\Z[x]/(g(x))$ and
$\Z[x]/(g^*(x))$ are non-isomorphic rings.
\end{lemma}

\begin{proof}
Let $\theta$ and $\theta^*$ denote the images of $x$ in $\Z[x]/(g(x))$
and $\Z[x]/(g^*(x))$, respectively.  By the assumption that $g$ and
$g^\ast$ are strongly quasi-reduced, we have that
$1,h_1(\theta),h_2(\theta),\ldots,h_{n-1}(\theta)$ and
$1,h_1^*(\theta^*),h_2^*(\theta^*),\ldots,h^*_{n-1}(\theta^*)$ are
the unique Minkowski-reduced bases of $\Z[x]/(g(x))$ and
$\Z[x]/(g^*(x))$, respectively, for some monic integer polynomials
$h_i$ and $h_i^*$ of degree $i$ for $i=1,\ldots,n-1$.

If $\phi:\Z[x]/(g(x))\to\Z[x]/(g^*(x))$ is a ring isomorphism, then by
the uniqueness of Minkowski-reduced bases for these rings, $\phi$ must
map Minkowski basis elements to Minkowski basis elements, i.e.,
$\phi(h_i(\theta))=h_i^*(\theta^*)$ for all $i$.  In particular, this
is true for $i=1$, so $\phi(\theta)=\theta^*+c$ for some $c\in\Z$,
since $h_1$ and $h_1^*$ are monic integer linear polynomials.
Therefore $\theta$ and $\theta^*+c$ must have the same minimal
polynomial, i.e., $g(x)=g^*(x-c)$; the assumption that $\theta$ and
$\theta^*$ both have trace 0 then implies that $c=0$.  It follows that
$g(x)=g^*(x)$, a contradiction.  We conclude that $\Z[x]/(g(x))$ and
$\Z[x]/(g^*(x))$ must be non-isomorphic rings, as desired.
\end{proof}

The condition of being quasi-reduced
%(and thus strictly quasi-reduced, since generic lattices have a unique Minkowski basis)
is fairly easy to attain:

\begin{lemma}
\source{squarefree-values-polynomial-discriminants}
\notready\label{quasilemma}
\source{squarefree-values-polynomial-discriminants}
If $g(x)$ is a monic real polynomial of nonzero discriminant, then
$g(\rho x)$ is quasi-reduced for any sufficiently large $\rho>0$.
\end{lemma}

\begin{proof}
This is easily seen from the Iwasawa-decomposition description of
Minkowski reduction. Consider the fundamental domain $\FF_\SL$ for
the action of $\SL_n(\Z)$ on $\SL_n(\R)$ given by
\begin{equation*}
\FF_\SL=\{\gamma=\nu \tau \kappa:\;\nu\in N'\;\tau\in T';\kappa\in\SO_n(\R)\},
\end{equation*}
where $N'$ denotes a compact subset (depending on $\tau$) of the group
of lower-triangular matrices and $T'$ is the group of diagonal
matrices $(t_1,\ldots,t_n)$ with $t_i\leq c\,t_{i+1}$ for all $i$ and
some absolute constant $c=c_n>0$.  Given an $n$-ary positive definite
integer-valued quadratic form $Q$, viewed as a symmetric $n\times n$
matrix, we write $Q=\gamma I_n \gamma^T$, where $I_n$ is the
sum-of-$n$-squares diagonal quadratic form and
$\gamma=\nu\tau\in\SL_n(\R)$ is unique up to right multiplication by
an element in $\SO_n(\R)$. The condition that $Q$ is Minkowski reduced
is equivalent to the condition that $\gamma$ belongs to $\FF_\SL$.
The condition that $Q$ be quasi-reduced is simply then that $t_i\leq
c\,t_{i+1}$ (with no condition on $\nu$).

%the condition that $Q$ is Minkowski reduced is equivalent to
%$Q=\gamma I_n \gamma^T$,
%where $I_n$ is the sum-of-$n$-squares
%diagonal quadratic form and $\gamma=\nu \tau \kappa$, where $\nu\in
%N'$, $\tau\in T'$, and $\kappa\in K$; here $N'$ as before denotes a
%compact subset (depending on $\tau$) of the group $N$ of
%lower-triangular matrices, $T'$ is the group of diagonal matrices
%$(t_1,\ldots,t_n)$ with $t_i\leq c\,t_{i+1}$ for all $i$ and some
%absolute constant $c=c_n>0$, and $K$ is the orthogonal group
%stabilizing the quadratic form $I_n$.


Consider the natural isomorphism $\R[x]/(g(x))\to \R[x]/(g(\rho x))$
of \'etale $\R$-algebras defined by $x\to\rho x$.  If $\theta$ denotes
the image of $x$ in $\R[x]/(g(x))$, then $\rho\theta$ is the image of
$x$ in $\R[x]/(g(\rho x))$ under this isomorphism.
Let $Q_1$ be the Gram matrix of the lattice basis
$1,\theta,\theta^2,\ldots,\theta^{n-1}$ in
$\R^{n}$ and $Q_\rho$ be the Gram matrix of the lattice basis
$1,\rho\theta,\rho^2\theta^2,\ldots,\rho^{n-1}\theta^{n-1}$ in
$\R^{n}$.  If the element $\tau \in T$
corresponding to $g(x)$ is $(t_1,\ldots,t_n)$, then the element
$\tau_\rho\in T$ corresponding to $g(\rho x)$ is $(t_1,\rho t_2,\rho^2
t_3,\ldots,\rho^{n-1}t_n)$. This is because $Q_\rho=\Lambda
Q_1\Lambda^T$, where $\Lambda$ is the diagonal matrix
$(1,\rho,\rho^2,\ldots,\rho^{n-1})$; therefore, if
$Q_1=(\nu\tau\kappa)I_n(\nu\tau\kappa)^T$, then $$Q_\rho=
(\Lambda\nu\tau\kappa)I_n(\Lambda\nu\tau\kappa)^T=(\nu'(\Lambda\tau)\kappa)I_n(\nu'(\Lambda\tau)\kappa)^T$$
for some $\nu'\in N$ depending on $\Lambda$, so
$\tau_\rho=\Lambda\tau$. For sufficiently large $\rho$, we then have
$\rho^{i-1}t_i\leq c\rho^{i}t_{i+1}$ for all $i=1,\ldots,n-1$, as
desired.
 \end{proof}

 \noindent
 Lemma~\ref{quasilemma} implies that most monic irreducible integer
 polynomials are strongly quasi-reduced:

\begin{lemma}
\source{squarefree-values-polynomial-discriminants}
\notready\label{mostqr}
\source{squarefree-values-polynomial-discriminants}
A density of $100\%$ of irreducible monic integer polynomials
$f(x)=x^n+a_1x^{n-1}+\cdots+a_n$ of degree~$n$, when ordered by
height $H(f):=\max\{|a_1|,|a_2|^{1/2},\ldots,|a_n|^{1/n}\}$,
are strongly quasi-reduced.
\end{lemma}

\begin{proof}
Let $\epsilon>0$, and let $B$ be a compact region in $\R^n\cong
\Vmn(\R)$ consisting of monic real polynomials of nonzero discriminant
and height less than $1$ such that $$\Vol(B)>(1-\epsilon)\Vol(\{f\in
\Vmn(\R):H(f)<1\}).$$ For each $f\in B$, by Lemma~\ref{quasilemma}
there exists a minimal finite constant $\rho_f>0$ such that $f(\rho
x)$ is quasi-reduced for any $\rho>\rho_f$. The function $\rho_f$ is
continuous in $f$, and thus by the compactness of $B$ there exists a
finite constant $\rho_B>0$ such that $f(\rho x)$ is quasi-reduced for
any $f\in B$ and $\rho>\rho_B$.

Now consider the weighted homogeneously expanding region $\rho\cdot B$
in $\R^n\cong \Vmn(\R)$, where a real number $\rho>0$ acts on $f\in B$
by $(\rho\cdot f)(x)=f(\rho x)$.  Note that $H(\rho\cdot f)=\rho
H(f)$.  For $\rho>\rho_B$, we have that all polynomials in $\rho\cdot
B$ are quasi-reduced, and
$$\Vol(\rho\cdot B)>(1-\epsilon)\Vol(\{f\in \Vmn(\R):H(f)<\rho\}).$$
Letting $\rho$ tend to infinity shows that the density of monic
integer polynomials $f$ of degree $n$, when ordered by height, that
have nonzero discriminant and are strongly quasi-reduced is greater
than $1-\epsilon$. Since $\epsilon$ was arbitrary, and $100\%$ of integer
polynomials are irreducible, the lemma follows.
\end{proof}

We have the following variation of Theorem~\ref{polydisc2}.

\begin{theorem}
\source{squarefree-values-polynomial-discriminants}
\notready\label{vanishinga1}
\source{squarefree-values-polynomial-discriminants}
Let $n\geq1$ be an integer.  Then when
monic integer polynomials $f(x)=x^n+a_1x^{n-1}+\cdots+a_n$ of
degree~$n$ with $a_1=0$ are ordered by $H(f):=
\max\{|a_1|,|a_2|^{1/2},\ldots,|a_n|^{1/n}\}$, the density having
squarefree discriminant $\Delta(f)$ exists and is equal to
$\kappa_n=\prod_p\kappa_n(p)>0$, where $\kappa_n(p)$ is the density of
monic polynomials $f(x)$ over $\Z_p$ with vanishing
$x^{n-1}$-coefficient having discriminant indivisible by $p^2$.
\end{theorem}

Indeed, the proof of Theorem~\ref{polydisc2} applies also to those
monic integer polynomials having vanishing $x^{n-1}$-coefficient %,
without any essential change; one simply replaces the representation
$W$ (along with $W_0$ and $W_{00}$) by the codimension-1 linear
subspace consisting of symmetric matrices with anti-trace $0$, but
otherwise the proof carries through in the identical manner.  The
analogue of Theorem~\ref{vanishinga1} holds also if the condition
$a_1=0$ is replaced by the condition $0\leq a_1<n$; in this case,
$\kappa_n=\prod_p\kappa_n(p)>0$ is replaced by the same constant
$\lambda_n=\prod_p\lambda_n(p)>0$ of Theorem~\ref{polydisc2}, since
for any monic degree-$n$ polynomial $f(x)$ there is a unique constant
$c\in\Z$ such that $f(x+c)$ has $x^{n-1}$-coefficient $a_1$ satisfying
$0\leq a_1<n$.


Lemmas \ref{reducedg} and \ref{mostqr} and Theorem~\ref{vanishinga1}
imply that 100\% of monic integer irreducible polynomials having
squarefree discriminant and vanishing $x^{n-1}$-coefficient (or those
having $x^{n-1}$-coefficient non-negative and less than $n$), when
ordered by height, yield {\it distinct} degree-$n$ fields.  Since
polynomials of height less than $X^{1/(n(n-1))}$ have absolute
discriminant $\ll X$, and since number fields of degree $n$ and
squarefree discriminant always have associated Galois group $S_n$, we
see that the number of $S_n$-number fields of degree $n$ and absolute
discriminant less than $X$ is $\gg
X^{(2+3+\cdots+n)/(n(n-1))}=X^{1/2+1/n}$. We have proven
Corollary~\ref{monogenic}.

\begin{remark}
\source{squarefree-values-polynomial-discriminants}
\notready{\em
The statement of Corollary~\ref{monogenic} holds even if one specifies
the real signatures of the monogenic $S_n$-number fields of degree
$n$, with the identical proof.  It holds also if one imposes any
desired set of local conditions on the degree-$n$ number fields at a
finite set of primes, so long as these local conditions do not
contradict local monogeneity.  }\end{remark}

\begin{remark}
\source{squarefree-values-polynomial-discriminants}
\notready{\em
We conjecture that a positive proportion of monic integer polynomials
of degree~$n$ with $x^{n-1}$-coefficient non-negative and less than
$n$ and absolute discriminant less than $X$ have height $O(
X^{1/(n(n-1))})$, where the implied $O$-constant depends only on $n$.
That is why we conjecture that the lower bound in
Corollary~\ref{monogenic} also gives the correct order of magnitude
for the upper bound.

In fact, let $C_n$ denote the $(n-1)$-dimensional Euclidean volume of
the $(n-1)$-dimensional region $R_0$ in $\Vmn(\R)\cong\R^n$
consisting of all polynomials $f(x)$ with vanishing
$x^{n-1}$-coefficient and absolute discriminant less than 1.  Then the
region $R_z$ in $\Vmn(\R)\cong\R^n$ of all polynomials $f(x)$ with
$x^{n-1}$-coefficient equal to $z$ and absolute discriminant less than
1 also has volume $C_n$, since $R_z$ is obtained from $R_0$ via the
volume-preserving transformation $x\mapsto x+z/n$.  Since we expect
that 100\% of monogenic number fields of degree $n$ can be expressed
as $\Z[\theta]$ in exactly one way (up to transformations of the form
$\theta\mapsto \pm \theta + c$ for $c\in \Z$), in view of
Theorem~\ref{polydiscmax2} we conjecture that the number of monogenic
number fields of degree $n$ and absolute discriminant less than $X$ is
asymptotic to
\begin{equation}
\frac{nC_n}{2\zeta(2)} X^{1/2+1/n}.
\end{equation}
%i.e., $1/\zeta(2)$ times the sum of the volumes of $R_0$, $R_1$, $\ldots$, $R_{n-1}$.
When $n=3$, a Mathematica computation shows that we have
$C_3= \frac{2^{1/3}(3+\sqrt{3})}{45} \frac{\Gamma(1/2)\Gamma(1/6)}{\Gamma(2/3)}$}.
\end{remark}

Finally, we turn to the proof of
Corollary~\ref{shortvector}. Following \cite{EV}, for any algebraic
number $x$, we write $\|x\|$ for the maximum of the archimedean
absolute values of~$x$. Given a number field~$K$, write
$s(K)=\inf\{\|x\|: x\in\OO_K,\; \Q(x)=K\}$. We consider the number of
number fields $K$ of degree~$n$ such that $s(K)\leq Y$.

As already pointed out in \cite[Remark~3.3]{EV}, an upper bound of $\ll
Y^{(n-1)(n+2)/2}$ is easy to obtain.  Namely, a bound on the
archimedean absolute values of an algebraic number~$x$ gives a bound
on the archimedean absolute values of all the conjugates of $x$, which
then gives a bound on the coefficients of the minimal polynomial of
$x$. Counting the number of possible minimal polynomials satisfying
these coefficient bounds gives the desired upper bound.

To obtain a lower bound of $\gg Y^{(n-1)(n+2)/2}$, we use
Lemmas~\ref{reducedg} and \ref{mostqr} and
Theorem~\ref{vanishinga1}. Suppose $f(x)=x^n + a_2x^{n-2} + \cdots +
a_n$ is an irreducible monic integer polynomial of degree $n$.
Let~$\theta$ denote a root of $f(x)$.
If $H(f)\leq Y$, then $|\theta|\ll Y$;
%Let $h(\theta)$ denote the height of $\theta$ as an algebraic number.
%Let $\theta_1,\ldots,\theta_n$ denote the conjugates of $\theta$ and
%for each $i=1,\ldots,n$, let $S_i$ denote the $i$-th elementary
%symmetric polynomial in $\theta_1,\ldots,\theta_n$. Suppose $H(f)\ll
%Y$. Then each $|S_i|\ll Y^i$. We claim that this implies that each
%$|\theta_i|\ll Y$ and so $\|\theta\|\ll Y$. Since $S_n\ll Y^n$, we see
%that after renaming, $|\theta_1|\ll Y$. For each $i=1,\ldots,n-1$, let
%$S_i'$ denote the $i$-th elementary symmetric polynomial in
%$\theta_2,\ldots,\theta_n$ and set $S_0'=1$. Then we have for every
%$i=1,\ldots,n-1$, $$S_i' = S_i - \theta_1S_{i-1}'.$$ Hence $|S_i'|\ll
%Y^i$ for every $i=1,\ldots,n-1$ and so by induction, $|\theta_i|\ll Y$
%for every $i=2,\ldots,n$.
this follows, e.g., from Fujiwara's bound~\cite{Fujiwara}:
$$ \|\theta\|\leq  \max\{ |a_1|,|a_2|^{1/2},\ldots,|a_{n-1}|^{1/(n-1)}|, |a_n/2|^{1/n}\}.$$

Therefore, if $H(f)\leq Y$, then
\begin{equation}\label{eq:est}
\source{squarefree-values-polynomial-discriminants}
s(\Q[x]/(f(x))) \leq \|\theta\| \ll Y.
\end{equation}
Now Lemma~\ref{mostqr} and Theorem \ref{vanishinga1} imply that there are
$\gg Y^{(n-1)(n+2)/2}$ such polynomials $f(x)$ of height less than $Y$ that have squarefree discriminant and are also strongly
quasi-reduced. Lemma~\ref{reducedg} and \eqref{eq:est} then imply that
these polynomials define distinct $S_n$-number fields $K$ of degree $n$ with
$s(K)\leq Y$. This completes the proof of Corollary~\ref{shortvector}.



\subsection*{Acknowledgments}

We thank Levent Alpoge, Benedict Gross, Wei Ho, Kiran Kedlaya, Hendrik
Lenstra, Barry Mazur, Bjorn Poonen, Peter Sarnak, and Ila Varma for
their kind interest and many helpful conversations.  The first and third
authors were supported by a Simons Investigator Grant and NSF Grant DMS-1001828.

\bigskip


\begin{thebibliography}{32}

\bibitem{ABZ}
A.\ Ash, J.\ Brakenhoff, and T.\ Zarrabi, Equality of Polynomial and Field Discriminants,
{\it Experiment.\ Math.} {\bf 16} (2007), 367--374.


\bibitem{geosieve} M.\ Bhargava, The geometric sieve and the density
  of squarefree values of polynomial discriminants and other invariant
  polynomials, {\tt http://arxiv.org/abs/1402.0031}.

\bibitem{AITBG} M.\ Bhargava and B.\ Gross, Arithmetic invariant theory,
  {\it Symmetry: representation theory and its applications,
    Progr. Math.}, {\bf 257}, Birkhuser/Springer, New York, (2014), 33--54.

\bibitem{BG2} M.\ Bhargava and B.\ Gross, The average size of the
  2-Selmer group of Jacobians of hyperelliptic curves having a
  rational Weierstrass point, {\it Automorphic representations and
    L-functions}, 23--91, Tata Inst.\ Fundam.\ Res.\ Stud.\ Math.\,
  \textbf{22}, Mumbai, 2013.


\bibitem{BGWhyper}
M.\ Bhargava, B.\ Gross. and X.\ Wang, A positive proportion of
locally soluble hyperelliptic curves over $\Q$ have no point over
any odd degree extension, {\it J. Amer. Math. Soc.} {\bf 30} (2017),
451--493.


\bibitem{SqSieve2} M.\ Bhargava, A.\ Shankar, and X.\ Wang, Squarefree
  values of polynomial discriminants II, in preperation.


\bibitem{Dietmann} R.\ Dietmann, On the distribution of Galois groups,
  {\it Mathematika} {\bf 58} (2012), 35--44.

\bibitem{Ek}
T.\ Ekedahl, An infinite version of the Chinese remainder theorem,
{\it Comment.\ Math.\ Univ.\ St.\ Paul.} {\bf 40} (1991), 53--59.

\bibitem{EV} J.\ Ellenberg and A.\ Venkatesh, The number of extensions
  of a number field with fixed degree and bounded discriminant, {\it
    Ann.\ of Math.} (2) {\bf 163} (2006), no.\ 2, 723--741.


\bibitem{Fujiwara} M.\ Fujiwara, {\"U}ber die obere Schranke des
  absoluten Betrages der Wurzeln einer algebraischen Gleichung, {\it
    Tohoku Math.\ Journal, First Series}, {\bf 10} (1916), 167--171.

\bibitem{GKZ} I.\ M.\ Gelfond, M.\ M.\ Kapranov and
  A.\ V.\ Zelevinsky, {\it Discriminants, Resultants and
    Multidimensional Determinants}, Springer Science + Business Media,
  New York, 1994.

\bibitem{Gabc} A.\ Granville, ABC allows us to count squarefrees,
  {\it Int.\ Math.\ Res.\ Notices.}\ (1998), no.\ 19, 991--1009


\bibitem{IK} H.\ Iwaniec and E.\ Kowalski, {\it Analytic Number Theory},
  {Amer.\ Math.\ Soc.\ Colloq.\ Publ.} {\bf 53}, Providence, RI (2004).

\bibitem{Kedlaya}
K.\ S.\ Kedlaya, A construction of polynomials with squarefree discriminants,
{\it Proc.\ Amer.\ Math.\ Soc.} {\bf 140} (2012), 3025--3033.

\bibitem{Kondo}
T.\ Kondo, Algebraic number fields with the discriminant equal to that of a
quadratic number field, {\it J.\ Math.\ Soc.\ Japan} {\bf 47} (1995),
no.\ 1, 31--36.

\bibitem{Nakagawa}
J.\ Nakagawa, On the Galois group of a number field with square free discriminant,
{\it Comment.\  Math.\  Univ.\ St.\ Paul.} {\bf 37} (1988), 95--98.

\bibitem{PBertini} B.\ Poonen, Bertini theorems over finite fields,
  {\it Ann.\ of Math}.\ (2004), 1099--1127.

\bibitem{PS2}
B.\ Poonen and M.\ Stoll,
Most odd degree hyperelliptic curves have only one rational point,
{\it Ann.\ of Math.\ }(2) \textbf{180} (2014), no.\ 3, 1137-–1166.

\bibitem{SK}
M.\ Sato and T.\ Kimura, A classification of irreducible
prehomogeneous vector spaces and their relative invariants,
{\it Nagoya Math.\ J.} {\bf 65} (1977), 1--155.

\bibitem{Scholz}
A.\ Scholz,  Aufgabe $208$,
{\it Jahresber.\ Deutsch Math.-Verein.}~{\bf XLV} (1935), 2.\ Abt.\ Heft 9/11, 110;
L\"osung der Aufgabe $208$,  {\it Jahresber.\ Deutsch Math.-Verein}.~{\bf XLVII} (1937), 2.\ Abt.\ Heft 5/8, 47.

\bibitem{SST} A.\ Shankar, A.\ S{\"o}dergren, and N.\ Templier,
  Sato-Tate equidistribution of certain families of Artin
  $L$-functions, {\tt https://arxiv.org/abs/1507.07031}.

\bibitem{ShTs} A.\ Shankar and J.\ Tsimerman, Counting $S_5$-fields with a
  power saving error term, {\it Forum Math.\ Sigma} \textbf{2} (2014), e13 (8 pages).

\bibitem{SW} A.\ Shankar and X.\ Wang, Rational points on hyperelliptic curves having a marked non-Weierstrass point. {\it Compositio Mathematica} \textbf{154} (2018). no.\ 1, 188--222.

\bibitem{W} X.\ Wang, {\it Pencils of quadrics and Jacobians of
  hyperelliptic curves}, Ph.D thesis, Harvard, 2013.

\bibitem{Yamamura}
K.\ Yamamura,
On unramified Galois extensions of real quadratic number fields,
{\it Osaka J.\ Math.} {\bf 23} (1986), no.\ 2, 471--478.

\bibitem{Yamamura2}
K.\ Yamamura, Some analogue of Hilbert's irreducibility theorem and
the distribution of algebraic number fields, {\it J.\ Fac.\ Sci.\ Univ.\ Tokyo Sect.\ IA Math.}
{\bf 38} (1991), no.\ 1, 99--135.



\end{thebibliography}
